\section{Additional work}
While researching the minimax and alpha-beta pruning algorithms, we found two interesting and promising alternatives/additions that we plan to explore in the coming week.\\
First, the Monte Carlo Tree Search (MCTS). This algorithm can be used to build an asymmetric decision tree to help explore certain promising areas of the decision tree. The idea is to explore areas of the tree and then (when finding a promising path) extending the tree with new nodes.\\
What is interesting is that because of its randomness, it is able to make meaningful evaluations without any heuristics. This means that when using our heuristics, the algorithm is less dependent on them and not at the mercy of their usability in a given instance. More importantly, however, is most likely the better \textit{anytime behaviour}. Because of the strict time limit, our computer player has to be able to stop its calculation suddenly and return the best known move. As the strength of MCTS grows 'smoothly' relative to the number of iterations/calculation time, we get a higher guarantee for "good" results.\\
This ties-in with our second addition: Iterative deepening search (IDS). The reason for better \textit{anytime behaviour} of MCTS compared to minimax (with or without alpha-beta pruning) is, that minimax is a depth-first algorithm. A depth-first algorithm aims to search the entire decision tree, however, because of the limited time span, it can be advantageous to only explore certain promising parts more deeply (like the MCTS does).\\
Because a DFS would miss nearby solutions by tending to explore long paths, we would want to use a BFS instead. However, because of the high branching factor, the BFS would most likely use too much memory. To solve this problem, IDS performs a DFS with limited depth repeatedly while increasing the depth limit. Hereby, it visits all nodes down to a certain depth d before visiting nodes in depth (d + 1), thereby behaving like a BFS, however using far less memory (as only DFS are repeatedly performed).
This would potentially allow us to reach better \textit{anytime behaviour} with minimax as well.